\chapter{Finite Elements, Spectral Methods, and Integral Operators}\label{ch:9}
\begin{chaptersummary}
The Galerkin framework connects error analysis to space selection and testing. Local finite-element assembly leads to global Fourier, Legendre, and Chebyshev representations, followed by integral-operator discretization with MoM. Support determines coupling range, while regularity and quadrature govern error, linking approximation theory to matrix structure and solvers in Volume 0.
\end{chaptersummary}
\learninggoals{Assemble one-dimensional linear finite-element mass and stiffness matrices; distinguish spectral Galerkin, collocation, and pseudospectral methods; construct a MoM matrix from an integral kernel; explain how smoothness controls approximation rates.}
\studyroute{Higher-dimensional element families, general spectral-error proofs, singular integration, and full radiation-condition treatment are omitted. Matrix solvers and fast operator application are developed in Volume~0.}{Galerkin finite element; spectral collocation; Fourier pseudospectral; Chebyshev approximation; method of moments}{For one-dimensional Poisson, specify the degrees of freedom, boundary constraints, and matrix entries for local linear FEM and global Legendre Galerkin approximation.}
\section{Local Finite-Element Spaces and Assembly}\index{Finite element method}
Finite elements start from a weak form. For Poisson's equation,
\[
-\nabla^2u=f
\]
multiply by a test function and integrate by parts:
\[
\int_\Omega\nabla v\cdot\nabla u\dd V
=
\int_\Omega vf\dd V+\text{boundary terms}.
\]

\begin{askai}
Ask AI:
\begin{quote}
Discretize one-dimensional Poisson using second-order central differences and linear finite elements. Write both matrices for five interior nodes and compare how differential-operator locality becomes matrix sparsity.
\end{quote}

Follow-up:
\begin{quote}
Using the same Poisson problem, explain the relationships among strong form, weak form, test function, basis function, and stiffness matrix.
\end{quote}
\end{askai}

Expanding
\[
u_h=\sum_jU_j\phi_j
\]
gives
\[
\bm K\bm U=\bm F.
\]

The expansion turns the weak form into equations for basis coefficients. A one-dimensional linear element makes the local matrices explicit, after which shared nodes determine assembly.


On a one-dimensional element $[x_j,x_{j+1}]$ of length $h_j$, the linear shape functions are
$\phi_0=1-\xi$, $\phi_1=\xi$, $\xi=(x-x_j)/h_j$.
Direct integration gives
\[
K^{(j)}=\frac1{h_j}\begin{pmatrix}1&-1\\-1&1\end{pmatrix},
\qquad
M^{(j)}=\frac{h_j}6\begin{pmatrix}2&1\\1&2\end{pmatrix}.
\]
Shared nodes represent the same degree of freedom; local entries are added using the node map.
Conforming spaces satisfy the required trace conditions. Scalar $H^1$ elements serve Poisson problems;
Maxwell curl problems require the corresponding tangential-trace structure, developed in Volume~2.

\conceptbridge{How can the same global polynomial space produce both Galerkin and collocation methods?}
{Modal coefficients and nodal values are representation choices. Galerkin constrains residuals by integral pairing; collocation constrains differential residuals at selected nodes. Suitable quadrature may relate the resulting matrices, with the precise relationship determined by testing, integration, and boundary treatment.}
{Global approximation space $\longrightarrow$ modal or nodal degrees of freedom $\longrightarrow$ residual mechanism.}

\section{Global Modes, Collocation, and Pseudospectral Methods}\index{Spectral methods}
Local finite elements organize assembly through compact support; global spectral representations exploit overall smoothness for accuracy. On the periodic interval $[0,2\pi)$, Fourier Galerkin takes
\[
u_N(x)=\sum_{k=-N}^{N}\widehat u_k\ee^{\ii kx},\qquad
\widehat{(\partial_xu_N)}_k=\ii k\widehat u_k.
\]
Galerkin tests against the retained Fourier modes. Pseudospectral evaluation transforms modes to nodal values,
computes nonlinear products in physical space, and transforms back. Unresolved high wave numbers may alias into lower ones;
appropriate padding or truncation removes aliasing for quadratic nonlinearities.

On $[-1,1]$, Legendre polynomials satisfy
\[
\int_{-1}^1P_jP_k\dd x=\frac{2}{2k+1}\delta_{jk}.
\]
Chebyshev polynomials $T_k(\cos\theta)=\cos(k\theta)$
are orthogonal with weight $(1-x^2)^{-1/2}$. This weighted inner product must be distinguished from the unweighted physical $L^2$ inner product.
Spectral collocation can use Chebyshev--Lobatto nodes $x_j=\cos(j\pi/p)$,
with $D_{ij}=\ell_j'(x_i)$. Differential residuals are constrained at interior nodes, and boundary values at boundary nodes.

For sufficiently smooth solutions with suitable analytic continuation, increasing the global degree can yield rapid or exponential convergence.
Finite regularity or discontinuities generally yield algebraic convergence, with possible Gibbs oscillations near interfaces.
Local polynomials facilitate geometry adaptation and refinement; global modes exploit global smoothness.
Error comparisons should report total degrees of freedom, quadrature error, and temporal error.

\section{Integral Operators and the Method of Moments}\index{Method of moments}
The method of moments in computational electromagnetics typically starts from an integral equation:
\[
\Lop J=g
\]
Expand
\[
J=\sum_{n=1}^N I_nf_n,
\]

\begin{authornote}
\textbf{Potential Question:}
"Why switch from differential equations to an integral equation and a current expansion?"

\textbf{Author Supplement:}
For certain linear scattering or radiation problems, a known background Green's function represents the field through equivalent boundary sources or currents. The unknown changes from a volume field to a boundary density determined by boundary conditions. This can reduce dimension, but the long-range Green's function generally produces a dense matrix.

\textbf{Core Conceptual Connection:}
\[
\text{field PDE}
\longrightarrow
\text{integral equation}
\longrightarrow
\text{current expansion}
\]
\end{authornote}

and constrain the integral-equation residual with test functions \(w_m\):
\[
\avg{w_m,\Lop J}
=
\avg{w_m,g}.
\]

\begin{askai}
Ask AI:
\begin{quote}
Discretize a simple one-dimensional integral equation using three piecewise-constant basis functions:
\[
\Lop J=g
\]
Write the entries \(\avg{w_m,\Lop f_n}\) explicitly and explain why each test function generally couples to every basis function.
\end{quote}

Follow-up:
\begin{quote}
Compare the nonzero patterns of differential- and integral-equation matrices, and identify the costs targeted by fast multipole methods or hierarchical matrices.
\end{quote}
\end{askai}

This gives
\[
\bm Z\bm I=\bm V.
\]

The long-range kernel generally makes $\bm Z$ dense. This structure and fast matrix application lead to Volume~0.

Specifically, if $(\mathcal LJ)(x)=\int_\Gamma G(x,y)J(y)\dd S_y$,
then
\[
Z_{mn}=\int_\Gamma w_m(x)\int_\Gamma G(x,y)f_n(y)\dd S_y\dd S_x,\qquad
V_m=\int_\Gamma w_mg\dd S.
\]
This testing pairing belongs to integral-equation discretization. Stability depends on the operator, trace spaces, kernel, and basis.
